\chapter{Viscosity Solutions}

\begin{center}
CHAPTER 5

\vspace{1em}

\Large \textbf{Viscosity Solutions}
\end{center}

In this chapter, we generalize the notion of classical solutions to viscosity solutions and study their regularities. We define viscosity solutions by comparing them with quadratic polynomials and thus remove the requirement that solutions be at least $C^2$. The main tool to study viscosity solutions is the maximum principle due to Alexandroff. We first generalize such a maximum principle to viscosity solutions and then use the resulting estimate to discuss the regularity of viscosity solutions. We use it to control the distribution functions of solutions and obtain Harnack inequality, which generalizes a result by Krylov and Safonov, and hence $C^\alpha$ regularity. We also use it to approximate solutions in $L^\infty$ by quadratic polynomials and get Schauder ($C^{2,\alpha}$) estimates. The methods are basically nonlinear, in the sense that they do not rely on differentiating equations. The results obtained in this chapter hold for general fully nonlinear equations, although in this note we focus only on linear equations.

\vspace{1em}

\begin{center}
\textbf{5.1. Alexandroff Maximum Principle}
\end{center}

We begin this section with the definition of viscosity solutions. This very weak concept of solutions enables us to define a class of functions containing all classical solutions of linear and nonlinear elliptic equations with fixed ellipticity constants and bounded measurable coefficients.

Suppose that $\Omega$ is a bounded and connected domain in $\mathbb{R}^n$ and that $a_{ij}\in C(\Omega)$ satisfies

\[
\lambda |\xi|^2 \le a_{ij}(x)\xi_i\xi_j \le \Lambda |\xi|^2 \qquad \text{for any } x\in \Omega \text{ and any } \xi\in\mathbb{R}^n,
\]

for some positive constants $\lambda$ and $\Lambda$. Consider the operator $L$ in $\Omega$ defined by

\[
Lu \equiv a_{ij}(x)D_{ij}u,
\]

for $u\in C^2(\Omega)$.

Suppose $u\in C^2(\Omega)$ is a supersolution in $\Omega$, i.e., $Lu\le 0$. Then for any $\varphi\in C^2(\Omega)$ with $L\varphi>0$, we have

\[
L(u-\varphi) < 0 \qquad \text{in } \Omega.
\]

This implies by the maximum principle that $u-\varphi$ cannot have local interior minimums in $\Omega$. In other words, if $u-\varphi$ has a local minimum at $x_0\in\Omega$, then

\[
L\varphi(x_0)\le 0.
\]

Geometrically, $u-\varphi$ having a local minimum at $x_0$ means that $\varphi$ touches $u$ from below at $x_0$ if we adjust $\varphi$ appropriately by adding a constant. This suggests the following definition. Here, we always assume $f\in C(\Omega)$.



\begin{definition}\label{hanlin:ch05-def-5-1}\dcref{Definition 5.1}\group{ch05-viscosity-solutions}\source{han-lin-lecture-notes:page-0102} $u \in C(\Omega)$ is a viscosity supersolution (resp. subsolution) of
$Lu = f$ in $\Omega$ if, for any $x_0 \in \Omega$ and any function $\varphi \in C^2(\Omega)$ such that $u - \varphi$ has a
local minimum (resp. maximum) at $x_0$, there holds
\[
L\varphi(x_0) \le f(x_0) \qquad (\text{resp. } L\varphi(x_0) \ge f(x_0)).
\]
We say that $u$ is a viscosity solution if it is a viscosity subsolution and a viscosity
supersolution.
\end{definition}

\begin{remark}\label{hanlin:ch05-rem-5-2}\dcref{Remark 5.2}\group{ch05-viscosity-solutions}\source{han-lin-lecture-notes:page-0102}\uses{hanlin:ch05-def-5-1} By an approximation, we may replace the $C^2$ function $\varphi$ by a
quadratic polynomial $Q$.
\end{remark}

\begin{remark}\label{hanlin:ch05-rem-5-3}\dcref{Remark 5.3}\group{ch05-viscosity-solutions}\source{han-lin-lecture-notes:page-0102}\uses{hanlin:ch05-def-5-1} The above analysis shows that a classical supersolution is a viscosity
supersolution. It is straightforward to prove that a $C^2$ viscosity supersolution
is a classical supersolution. Similar statements hold for subsolutions and solutions.

\end{remark}
\medskip


\noindent The notion of viscosity solutions can be generalized to nonlinear elliptic equations
accordingly.


\noindent Now we define in a weak way the class of “all solutions to all elliptic equations”.
For any function $\varphi$, which is $C^2$ at $x_0$, we have the following equivalence
\[
\sum_{i,j=1}^n a_{ij}(x_0)D_{ij}\varphi(x_0) \le 0
\]
\[
\Longleftrightarrow \sum_{k=1}^n \alpha_k e_k \le 0 \text{ with } \lambda \le \alpha_k \le \Lambda,\, e_k = e_k(D^2\varphi(x_0))
\]
\[
\Longleftrightarrow \sum_{e_i>0} \alpha_i e_i + \sum_{e_i<0} \alpha_i e_i \le 0
\]
\[
\Longleftrightarrow \sum_{e_i>0} \alpha_i e_i \le \sum_{e_i<0} \alpha_i(-e_i),
\]
which implies
\[
\lambda \sum_{e_i>0} e_i \le \Lambda \sum_{e_i<0} (-e_i),
\]
where $e_1,\cdots,e_n$ are eigenvalues of the Hessian matrix $D^2\varphi(x_0)$. This means that
positive eigenvalues of $D^2\varphi(x_0)$ are controlled by negative eigenvalues.

\begin{definition}\label{hanlin:ch05-def-5-4}\dcref{Definition 5.4}\group{ch05-viscosity-solutions}\source{han-lin-lecture-notes:page-0102} Suppose $f$ is a continuous function in $\Omega$ and that $\lambda$ and $\Lambda$
are two positive constants. We define $u \in C(\Omega)$ to belong to $S^+(\lambda,\Lambda,f)$ (resp.
$S^-(\lambda,\Lambda,f)$) if, for any $x_0 \in \Omega$ and any function $\varphi \in C^2(\Omega)$ such that $u - \varphi$ has a
local minimum (resp. maximum) at $x_0$, there holds
\[
\lambda \sum_{e_i>0} e_i(x_0) + \Lambda \sum_{e_i<0} e_i(x_0) \le f(x_0)
\]
\[
(\text{resp. } \Lambda \sum_{e_i>0} e_i(x_0) + \lambda \sum_{e_i<0} e_i(x_0) \ge f(x_0)),
\]
where $e_1(x_0),\cdots,e_n(x_0)$ are eigenvalues of the Hessian matrix $D^2\varphi(x_0)$.


\noindent We set $S(\lambda,\Lambda,f)=S^+(\lambda,\Lambda,f)\cap S^-(\lambda,\Lambda,f)$.
\end{definition}




\begin{remark}\label{hanlin:ch05-rem-5-5}\dcref{Remark 5.5}\group{ch05-viscosity-solutions}\source{han-lin-lecture-notes:page-0103}\uses{hanlin:ch05-def-5-4} Any viscosity supersolutions of
\[
a_{ij}D_{ij}u=f \quad \text{in }\Omega
\]
belong to the class $S^+(\lambda,\Lambda,f)$, if
\[
\lambda |\xi|^2 \le a_{ij}(x)\xi_i\xi_j \le \Lambda |\xi|^2
\qquad \text{for any } x\in\Omega \text{ and any } \xi\in\mathbf{R}^n.
\]

The class $S^+(\lambda,\Lambda,f)$ and $S^-(\lambda,\Lambda,f)$ also include solutions of fully nonlinear equations. Among them are the Pucci's equations.

\end{remark}
\begin{example}\label{hanlin:ch05-ex-5-6}\dcref{Example 5.6}\group{ch05-viscosity-solutions}\source{han-lin-lecture-notes:page-0103}\uses{hanlin:ch05-def-5-4} For any two positive constants $\lambda \le \Lambda$, let $A$ be a symmetric matrix whose eigenvalues belong to $[\lambda,\Lambda]$, i.e.,
\[
\lambda |\xi|^2 \le A_{ij}\xi_i\xi_j \le \Lambda |\xi|^2
\qquad \text{for any } \xi\in\mathbf{R}^n.
\]

Denote by $\mathcal{A}_{\lambda,\Lambda}$ the class of all such matrices. For any symmetric matrix $M$, we define the Pucci's extremal operators
\[
\mathcal{M}^-(M)=\mathcal{M}^-(\lambda,\Lambda,M)=\inf_{A\in\mathcal{A}_{\lambda,\Lambda}} A_{ij}M_{ij},
\]
\[
\mathcal{M}^+(M)=\mathcal{M}^+(\lambda,\Lambda,M)=\sup_{A\in\mathcal{A}_{\lambda,\Lambda}} A_{ij}M_{ij}.
\]

Pucci's equations are given by
\[
\mathcal{M}^-(\lambda,\Lambda,M)=f, \qquad \mathcal{M}^+(\lambda,\Lambda,M)=g,
\]
for continuous functions $f$ and $g$ in $\Omega$. It is easy to see
\[
\mathcal{M}^-(\lambda,\Lambda,M)=\lambda\sum_{e_i>0} e_i+\Lambda\sum_{e_i<0} e_i,
\]
\[
\mathcal{M}^+(\lambda,\Lambda,M)=\Lambda\sum_{e_i>0} e_i+\lambda\sum_{e_i<0} e_i,
\]
where $e_1,\dots,e_n$ are eigenvalues of $M$. Therefore $u\in S^+(\lambda,\Lambda,f)$ if and only if
\[
\mathcal{M}^-(\lambda,\Lambda,D^2u)\le f
\]
in the viscosity sense, i.e., for any $\varphi\in C^2(\Omega)$ such that $u-\varphi$ has a local minimum at $x_0\in\Omega$ there holds
\[
\mathcal{M}^-(\lambda,\Lambda,D^2\varphi(x_0))\le f(x_0).
\]

By the definition of $\mathcal{M}^-$ and $\mathcal{M}^+$, it is easy to check for any two symmetric matrices $M$ and $N$
\[
\mathcal{M}^-(M)+\mathcal{M}^-(N)\le \mathcal{M}^-(M+N)\le \mathcal{M}^+(M)+\mathcal{M}^-(N)
\le \mathcal{M}^+(M+N)\le \mathcal{M}^+(M)+\mathcal{M}^+(N).
\]
This property will be needed in the next section.

\end{example}

Next, we derive the Alexandroff maximum principle for viscosity solutions. It has the role of energy inequalities for solutions of equations of divergence form. Let $v$ be a continuous function in an open convex set $\Omega$. The convex envelope of $v$ in $\Omega$ is defined by
\[
\Gamma(v)(x)=\sup_L\{L(x);\, L\le v \text{ in }\Omega,\, L \text{ is an affine function}\},
\]
for any $x\in\Omega$. It is easy to see that $\Gamma(v)$ is a convex function in $\Omega$. The set
\[
\{v=\Gamma(v)\}=\{x\in\Omega; v(x)=\Gamma(v)(x)\}
\]



is called the (lower) contact set of $v$. Points in the contact set are called contact points.

\medskip


\noindent The following is the classical version of the Alexandroff maximum principle.
We do not require that functions be solutions to elliptic equations. See Lemma
2.33 in Chapter 2.

\medskip

\begin{lemma}\label{hanlin:ch05-lem-5-7}\dcref{Lemma 5.7}\group{ch05-viscosity-solutions}\source{han-lin-lecture-notes:page-0104} \emph{Suppose $u$ is a $C^{1,1}$ function in $B_1$ with $u \ge 0$ on $\partial B_1$. Then}
\[
\sup_{B_1} u^- \le c\left(\int_{B_1\cap\{u=\Gamma_u\}} \det D^2u\right)^{\frac{1}{n}},
\]
\emph{where $\Gamma_u$ is the convex envelope of $-u^- = \min\{u,0\}$ and $c$ is a positive constant depending only on $n$.}
\end{lemma}

\medskip


\noindent Now we state a version for viscosity solutions.

\medskip

\begin{theorem}\label{hanlin:ch05-thm-5-8}\dcref{Theorem 5.8}\group{ch05-viscosity-solutions}\source{han-lin-lecture-notes:page-0104,han-lin-lecture-notes:page-0105}\uses{hanlin:ch05-def-5-4,hanlin:ch05-lem-5-7} \emph{Suppose $u$ belongs to $S^+(\lambda,\Lambda,f)$ in $B_1$ with $u \ge 0$ on $\partial B_1$ for some $f\in C(\Omega)$. Then}
\[
\sup_{B_1} u^- \le c\left(\int_{B_1\cap\{u=\Gamma_u\}} (f^+)^n\right)^{\frac{1}{n}},
\]
\emph{where $\Gamma_u$ is the convex envelope of $-u^- = \min\{u,0\}$ and $c$ is a positive constant depending only on $n,\lambda$ and $\Lambda$.}
\end{theorem}

\medskip

\textbf{Proof.} We prove that $\Gamma_u$ is a $C^{1,1}$ function in $B_1$ and that at any contact
point $x_0$
\[
(1)\hspace{7em} f(x_0)\ge 0,
\]
and
\[
(2)\hspace{4.5em} L(x)\le \Gamma_u(x)\le L(x)+C\{f(x_0)+\varepsilon(x)\}|x-x_0|^2
\]
for some affine function $L$ and any $x$ close to $x_0$, where $\varepsilon(x)\to 0$ as $x\to x_0$ and $C$
is a positive constant depending only on $n,\lambda$ and $\Lambda$. We then obtain by (2)
\[
\det D^2\Gamma_u(x)\le C(f(x))^n \qquad \text{for a.e. } x\in\{u=\Gamma_u\}.
\]

Now we apply Lemma 5.7 to $\Gamma_u$ to get the desired result.
Suppose $x_0$ is a contact point, i.e., $u(x_0)=\Gamma_u(x_0)$. We may assume $x_0=0$.
We also assume, by subtracting a supporting plane at $x_0=0$, that $u\ge 0$ in $B_1$ and
that $u(0)=0$.

In order to prove (1), we take $h(x)=-\varepsilon|x|^2/2$ in $B_1$. Obviously, $u-h$ has
a minimum at $0$. Note that the eigenvalues of $D^2h(0)$ are $-\varepsilon,\cdots,-\varepsilon$. By the
definition of $S^+(\lambda,\Lambda,f)$, we have
\[
-n\Lambda\varepsilon\le f(0).
\]

We get (1) by letting $\varepsilon\to 0$.
For (2), we prove
\[
0\le \Gamma_u(x)\le C(n,\lambda,\Lambda)\{f(0)+\varepsilon(x)\}|x|^2 \qquad \text{for any } x\in B_1,
\]
where $\varepsilon(x)\to 0$ as $x\to 0$. By setting $w=\Gamma_u$, we need to estimate for any small
$r>0$
\[
C_r=\frac{1}{r^2}\max_{B_r} w.
\]


Fix an $r>0$. By the convexity, $w$ attains its maximum in $\bar B_r$ at some point on the boundary, say, $(0,\cdots,0,r)$. The set $\{x\in B_1;w(x)\leq w(0,\cdots,0,r)\}$ is convex and contains $B_r$. It follows easily that

\[
w(x',r)\geq w(0,\cdots,0,r)=C_r r^2 \qquad \text{for any } x=(x',r)\in B_1.
\]

Take a positive number $N$ to be determined and set

\[
R_r=\{(x',x_n);|x'|\leq Nr, |x_n|\leq r\}.
\]

We construct a quadratic polynomial that touches $u$ from below in $R_r$ and curves upward very much. Set for some $b>0$

\[
h(x)=(x_n+r)^2-b|x'|^2.
\]

Then we have

\begin{align*}
\text{(i)}\ & \text{for } x_n=-r,\ h\leq 0;\\
\text{(ii)}\ & \text{for } |x'|=Nr,\ h\leq (4-bN^2)r^2\leq 0 \text{ if we take } b=4/N^2;\\
\text{(iii)}\ & \text{for } x_n=r,\ h=4r^2-b|x'|^2\leq 4r^2.
\end{align*}

Hence if we set

\[
\tilde h(x)=\frac{C_r}{4}h(x)=\frac{C_r}{4}\left\{(x_n+r)^2-\frac{4}{N^2}|x'|^2\right\},
\]

we obtain $\tilde h\leq w\leq u$ on $\partial R_r$ (since $w$ is the convex envelope of $u$) and $\tilde h(0)=C_r r^2/4>0=w(0)=u(0)$. By lowering $\tilde h$ appropriately, we conclude that $u-\tilde h$ has a local minimum somewhere inside $R_r$. Note that the eigenvalues of $D^2\tilde h$ are given by $C_r/2,-2C_r/N^2,\cdots,-2C_r/N^2$. Hence by the definition of $S^+(\lambda,\Lambda,f)$, we have

\[
\lambda \frac{C_r}{2}-2\Lambda(n-1)\frac{C_r}{N^2}\leq \max_{R_r} f.
\]

By choosing $N$ large, depending only on $n,\lambda$ and $\Lambda$, we obtain

\[
C_r\leq \frac{4}{\lambda}\max_{R_r} f,
\]

or

\[
\max_{B_r} w\leq \frac{4}{\lambda}r^2\max_{R_r} f.
\]

We finish the proof by noting $\max_{R_r}f\to f(0)$ as $r\to 0$. \hfill$\square$

We end this section with a simple consequence of the Calderon-Zygmund decomposition. We first recall some terminology. Let $Q_1$ be the unit cube. Cut it equally into $2^n$ cubes, which we take as the first generation. Do the same cutting for these small cubes to get the second generation. Continue this process. These cubes (from all generations) are called dyadic cubes. Any $(k+1)$-generation cube $Q$ arises from some $k$-generation cube $\tilde Q$, which is called the predecessor of $Q$.

\medskip

\begin{lemma}\label{hanlin:ch05-lem-5-9}\dcref{Lemma 5.9}\group{ch05-viscosity-solutions}\source{han-lin-lecture-notes:page-0105,han-lin-lecture-notes:page-0106}\uses{hanlin:weaksol-part1-cz} \textit{Let $A\subset B\subset Q_1$ be measurable sets such that}
\begin{itemize}
\item[(i)] $|A|<\delta$ for some $\delta\in(0,1)$;
\item[(ii)] for any dyadic cube $Q$, $|A\cap Q|\geq \delta|Q|$ implies $\tilde Q\subset B$ for the predecessor $\tilde Q$ of $Q$.
\end{itemize}
\textit{Then}

\[
|A|\leq \delta|B|.
\]
\end{lemma}


Proof. Apply the Calderon-Zygmund decomposition, Lemma 3.7 in Chapter 3, to $f=\chi_A$. We obtain, by the assumption (i), a sequence of dyadic cubes $\{Q^j\}$ such that

\[
A \subset \bigcup_j Q^j \text{ except for a set of measure zero,}
\]

\[
\delta \leq \frac{|A\cap Q^j|}{|Q^j|} < 2^n\delta,
\]

and

\[
\frac{|A\cap \Qtilde^j|}{|\Qtilde^j|} < \delta,
\]

for any predecessor $\Qtilde^j$ of $Q^j$. By the assumption (ii), we have $\Qtilde^j \subset B$ for each $j$.
Hence, we obtain

\[
A \subset \bigcup_j \Qtilde^j \subset B.
\]

We relabel $\{\Qtilde^j\}$ so that they are nonoverlapping. Therefore, we get

\[
|A| \leq \sum_i |A\cap \Qtilde^i| \leq \delta \sum_i |\Qtilde^i| \leq \delta |B|.
\]

This finishes the proof.

\medskip


\noindent\textbf{5.2. The Harnack Inequality}

The main result in this section is the following Harnack inequality.

\medskip

\begin{theorem}\label{hanlin:ch05-thm-5-10}\dcref{Theorem 5.10}\group{ch05-viscosity-solutions}\source{han-lin-lecture-notes:page-0106,han-lin-lecture-notes:page-0107}\uses{hanlin:ch05-lem-5-12} \textit{Suppose $u$ belongs to $\Sclass(\lambda,\Lambda,f)$ in $\B_1$ for some $f\in C(\B_1)$ with $u\geq 0$ in $\B_1$. Then}

\[
\sup_{\B_{1/2}} u \leq c\left\{\inf_{\B_{1/2}} u + \|f\|_{L^n(\B_1)}\right\},
\]

\textit{where $c$ is a positive constant depending only on $n$, $\lambda$ and $\Lambda$.}
\end{theorem}

The interior Hölder continuity of solutions is a direct consequence, whose proof is identical to that of Theorem 4.17 in Chapter 4.

\medskip

\begin{corollary}\label{hanlin:ch05-cor-5-11}\dcref{Corollary 5.11}\group{ch05-viscosity-solutions}\source{han-lin-lecture-notes:page-0106}\uses{hanlin:ch05-thm-5-10} \textit{Suppose $u$ belongs to $\Sclass(\lambda,\Lambda,f)$ in $\B_1$ for some $f\in C(\B_1)$. Then $u\in C^\alpha(\B_1)$, for some $\alpha\in (0,1)$ depending only on $n$, $\lambda$ and $\Lambda$. Moreover,}

\[
|u(x)-u(y)| \leq c|x-y|^\alpha \left\{\sup_{\B_1} |u| + \|f\|_{L^n(\B_1)}\right\} \qquad \text{for any } x,y\in \B_{1/2},
\]

\textit{where $c$ is a positive constant depending only on $n$, $\lambda$ and $\Lambda$.}
\end{corollary}

For convenience, we work in cubes instead of balls. We prove the following result.

\medskip

\begin{lemma}\label{hanlin:ch05-lem-5-12}\dcref{Lemma 5.12}\group{ch05-viscosity-solutions}\source{han-lin-lecture-notes:page-0106,han-lin-lecture-notes:page-0109,han-lin-lecture-notes:page-0110,han-lin-lecture-notes:page-0111}\uses{hanlin:ch05-lem-5-13,hanlin:ch05-lem-5-14} \textit{Suppose $u$ belongs to $\Sclass(\lambda,\Lambda,f)$ in $Q_{4\sqrt{n}}$ for some $f\in C(Q_{4\sqrt{n}})$, with $u\geq 0$ in $Q_{4\sqrt{n}}$. Then there exist two positive constants $\varepsilon_0$ and $C$, depending only on $n$, $\lambda$ and $\Lambda$, such that, if $\inf_{Q_{1/4}} u \leq 1$ and $\|f\|_{L^n(Q_{4\sqrt{n}})} \leq \varepsilon_0$, then $\sup_{Q_{1/4}} u \leq C$.}
\end{lemma}

Theorem 5.10 follows from Lemma 5.12 easily. For any $u \in \Splus(\lambda,\Lambda,f)$ in $Q_{4\sqrt{n}}$
with $u \geq 0$ in $Q_{4\sqrt{n}}$, consider for $\delta > 0$

\[
u_{\delta}=\frac{u}{\inf_{Q_{1/4}} u+\delta+\frac{1}{\varepsilon_{0}}\|f\|_{L^{n}(Q_{4\sqrt{n}})}}.
\]

Applying Lemma 5.12 to $u_{\delta}$ and then letting $\delta \to 0$, we get

\[
\sup_{Q_{1/4}} u \leq C\left\{\inf_{Q_{1/4}} u+\|f\|_{L^{n}(Q_{4\sqrt{n}})}\right\}.
\]

Then Theorem 5.10 follows by a standard covering argument.
Now we begin to prove Lemma 5.12. The following result plays a key role. It
asserts that if a solution is small somewhere in $Q_{3}$ then it is under control in a good
portion of $Q_{1}$.

\medskip

\begin{lemma}\label{hanlin:ch05-lem-5-13}\dcref{Lemma 5.13}\group{ch05-viscosity-solutions}\source{han-lin-lecture-notes:page-0107,han-lin-lecture-notes:page-0108}\uses{hanlin:ch05-def-5-4,hanlin:ch05-thm-5-8} \emph{Suppose $u$ belongs to $\Splus(\lambda,\Lambda,f)$ in $B_{2\sqrt{n}}$ for some $f \in C(B_{2\sqrt{n}})$.
Then there exist constants $\varepsilon_{0} > 0$, $\mu \in (0,1)$ and $M > 1$, depending only on $n$, $\lambda$
and $\Lambda$, such that if}

\[
\begin{array}{ll}
(1) & u \geq 0 \text{ in } B_{2\sqrt{n}},\\
& \inf_{Q_{3}} u \leq 1,\\
& \|f\|_{L^{n}(B_{2\sqrt{n}})} \leq \varepsilon_{0},
\end{array}
\]

\emph{then}

\[
|\{u \leq M\} \cap Q_{1}| > \mu.
\]
\end{lemma}

\medskip


\noindent\textbf{Proof.} We construct a function $g$, which is very concave outside $Q_{1}$, such
that the contact set occurs in $Q_{1}$ if we correct $u$ by $g$. In other words, we localize
where the contact occurs by choosing suitable functions.

Note $B_{1/4} \subset B_{1/2} \subset Q_{1} \subset Q_{3} \subset B_{2\sqrt{n}}$. Define $g$ in $B_{2\sqrt{n}}$ by

\[
g(x)=-M\left(1-\frac{|x|^{2}}{4n}\right)^{\beta},
\]

for large $\beta > 0$ to be determined and some $M > 0$. We choose $M$, according to $\beta$,
such that

\[
(2)\qquad g=0 \text{ on } \partial B_{2\sqrt{n}} \text{ and } g \leq -2 \text{ in } Q_{3}.
\]

Set $w=u+g$ in $B_{2\sqrt{n}}$. We prove by choosing $\beta$ large such that

\[
(3)\qquad w \in \Splus(\lambda,\Lambda,f) \quad \text{in } B_{2\sqrt{n}} \setminus Q_{1}.
\]

Suppose $\varphi$ is a quadratic polynomial such that $w-\varphi$ has a local minimum at
$x_{0} \in B_{2\sqrt{n}}$. Then $u-(\varphi-g)$ has a local minimum at $x_{0} \in B_{2\sqrt{n}}$. By the definition
of $\Splus(\lambda,\Lambda,f)$ and the Pucci's extremal operator $\Mminus$, we have

\[
\Mminus(\lambda,\Lambda,D^{2}\varphi(x_{0})-D^{2}g(x_{0})) \leq f(x_{0}),
\]

or

\[
\Mminus(\lambda,\Lambda,D^{2}\varphi(x_{0}))+\Mminus(\lambda,\Lambda,-D^{2}g(x_{0})) \leq f(x_{0}),
\]

where we used the property of $\Mminus$. We choose $\beta$ large such that

\[
\Mminus(\lambda,\Lambda,-D^{2}g(x_{0})) \geq 0 \qquad \text{for any } x_{0} \in B_{2\sqrt{n}} \setminus B_{1/4}.
\]


To see this, we need to calculate the Hessian matrix of $g$. Note
\[
D_{ij}g(x)=\frac{M}{2n}\beta\Bigl(1-\frac{|x|^2}{4n}\Bigr)^{\beta-1}\delta_{ij}
-\frac{M}{(2n)^2}\beta(\beta-1)\Bigl(1-\frac{|x|^2}{4n}\Bigr)^{\beta-2}x_i x_j.
\]
If we choose $x=(|x|,0,\dots,0)$, then the eigenvalues of $-D^2g(x)$ are given by
\[
\frac{M}{2n}\beta\Bigl(1-\frac{|x|^2}{4n}\Bigr)^{\beta-2}
\Bigl(\frac{2\beta-1}{4n}|x|^2-1\Bigr),
\]
with the multiplicity $1$, and
\[
-\frac{M}{2n}\beta\Bigl(1-\frac{|x|^2}{4n}\Bigr)^{\beta-1},
\]
with the multiplicity $n-1$. We choose $\beta$ large so that for $|x|\geq 1/4$ the first eigenvalue is positive and the rest negative, denoted by $e^+(x)$ and $e^-(x)$ respectively. Therefore for $|x|\geq 1/4$, we have
\[
\M^-(\lambda,\Lambda,-D^2g(x))=\lambda e^+(x)+(n-1)\Lambda e^-(x)
\]
\[
=\frac{M}{2n}\beta\Bigl(1-\frac{|x|^2}{4n}\Bigr)^{\beta-2}
\left\{\lambda\left(\frac{2\beta-1}{4n}|x|^2-1\right)-(n-1)\Lambda\left(1-\frac{|x|^2}{4n}\right)\right\}\geq 0,
\]
if we choose $\beta$ large, depending only on $n,\lambda$ and $\Lambda$. This finishes the proof of (3).
In fact, we obtain
\[
 w\in\Splus(\lambda,\Lambda,f+\eta) \quad \text{in } B_{2\sqrt n}
\]
for some $\eta\in C_0^\infty(Q_1)$ and $0\leq \eta\leq C(n,\lambda,\Lambda)$.
We now apply Theorem 5.8 to $w$ in $B_{2\sqrt n}$. Note that $\inf_{Q_3} w\leq -1$ and $w\geq 0$
on $\partial B_{2\sqrt n}$ by (1) and (2). We obtain
\[
1\leq C\left(\int_{B_{2\sqrt n}\cap\{w=\Gamma_w\}} (|f|+\eta)^n\right)^{\frac1n}
\]
\[
\leq C\|f\|_{L^n(B_{2\sqrt n})}+C|\{w=\Gamma_w\}\cap Q_1|^{\frac1n}.
\]
Choosing $\varepsilon_0$ small enough, we get
\[
\frac12\leq C|\{w=\Gamma_w\}\cap Q_1|^{\frac1n}\leq C|\{u\leq M\}\cap Q_1|^{\frac1n},
\]
since $w(x)=\Gamma_w(x)$ implies $w(x)\leq 0$ and hence $u(x)\leq -g(x)\leq M$. This finishes the proof.

Next, we prove the power decay of distribution functions.

\begin{lemma}\label{hanlin:ch05-lem-5-14}\dcref{Lemma 5.14}\group{ch05-viscosity-solutions}\source{han-lin-lecture-notes:page-0108,han-lin-lecture-notes:page-0109}\uses{hanlin:ch05-lem-5-9,hanlin:ch05-lem-5-13} \textit{Suppose $u$ belongs to $\Splus(\lambda,\Lambda,f)$ in $B_{2\sqrt n}$ for some $f\in C(B_{2\sqrt n})$.
Then there exist positive constants $\varepsilon_0$, $\varepsilon$ and $C$, depending only on $n,\lambda$ and $\Lambda$, such
that if}
\[
\tag{1}
\begin{aligned}
&u\geq 0 \text{ in } B_{2\sqrt n},\\
&\inf_{Q_3} u\leq 1,\\
&\|f\|_{L^n(B_{2\sqrt n})}\leq \varepsilon_0,
\end{aligned}
\]
\textit{then}
\[
|\{u\geq t\}\cap Q_1|\leq Ct^{-\varepsilon}\text{ for any } t>0.
\]
\end{lemma}


\textsc{Proof.} We prove that, under the assumption (1), there holds
\begin{equation}\tag{2}
\left|\{u>M^k\}\cap Q_1\right|\leq (1-\mu)^k \text{ for } k=1,2,\ldots,
\end{equation}
where $M$ and $\mu$ are as in Lemma 5.13.

For $k=1$, (2) is simply Lemma 5.13. Suppose now (2) holds for $k-1$. Set
\[
A=\{u>M^k\}\cap Q_1,\qquad B=\{u>M^{k-1}\}\cap Q_1.
\]
We use Lemma 5.9 to prove
\begin{equation}\tag{3}
|A|\leq (1-\mu)|B|.
\end{equation}
Clearly $A\subset B\subset Q_1$ and $|A|\leq |\{u>M\}\cap Q_1|\leq 1-\mu$ by Lemma 5.13. We claim that, if $Q=Q_r(x_0)$ is a cube in $Q_1$ such that
\begin{equation}\tag{4}
|A\cap Q|>(1-\mu)|Q|,
\end{equation}
then $\tilde{Q}\cap Q_1\subset B$ for $\tilde{Q}=Q_{3r}(x_0)$. We prove it by contradiction. If not, we take a $\tilde{x}\in\tilde{Q}$ such that $u(\tilde{x})\leq M^{k-1}$. Consider the transformation
\[
x=x_0+ry \qquad \text{for any } y\in Q_1 \text{ and } x\in Q=Q_r(x_0),
\]
and the function
\[
\tilde{u}(y)=\frac{1}{M^{k-1}}u(x).
\]
Then $\tilde{u}\geq 0$ in $B_{2\sqrt{n}}$ and $\inf_{Q_3}\tilde{u}\leq 1$. It is easy to check that $\tilde{u}\in S^+(\lambda,\Lambda,\tilde{f})$ in $B_{2\sqrt{n}}$ with $\|\tilde{f}\|_{L^n(B_{2\sqrt{n}})}\leq \varepsilon_0$. In fact, we have
\[
\tilde{f}(y)=\frac{r^2}{M^{k-1}}f(x) \qquad \text{for any } y\in B_{2\sqrt{n}},
\]
and hence
\[
\|\tilde{f}\|_{L^n(B_{2\sqrt{n}})}
\leq \frac{r}{M^{k-1}}\|f\|_{L^n(B_{2\sqrt{n}})}
\leq \|f\|_{L^n(B_{2\sqrt{n}})}\leq \varepsilon_0.
\]
Then $\tilde{u}$ satisfies the assumption (1). We now apply Lemma 5.13 to $\tilde{u}$ to get
\[
\mu<\left|\{\tilde{u}(y)\leq M\}\cap Q_1\right|=r^{-n}\left|\{u(x)\leq M^k\}\cap Q\right|.
\]
Hence $|Q\cap A^c|>\mu|Q|$, which contradicts (4). We are in a position to apply Lemma 5.9 to get (3).
\hfill$\square$

\textsc{Proof of Lemma 5.12.} We prove that there exist two constants $\theta>1$ and $M_0\gg 1$, depending only on $n,\lambda$ and $\Lambda$, such that, if $u(x_0)=P>M_0$ for some $x_0\in B_{1/4}$, there exists a sequence $\{x_k\}\subset B_{1/2}$ such that
\[
u(x_k)\geq \theta^k P \qquad \text{for } k=0,1,2,\ldots.
\]
This contradicts the boundedness of $u$, and hence we conclude that $\sup_{B_{1/4}}u\leq M_0$.

Suppose $u(x_0)=P>M_0$ for an $x_0\in B_{1/4}$, with $M_0$ and $\theta$ to be determined in the process. Consider a cube $Q_r(x_0)$, centered at $x_0$ with the side length $r$, which will be chosen later. We intend to find a point $x_1\in Q_{4\sqrt{n}r}(x_0)$ such that $u(x_1)\geq \theta P$. To achieve this, we first choose $r$ such that $\{u>P/2\}$ covers less than half of $Q_r(x_0)$. This can be done by using the power decay of the distribution function of $u$.

Note $\inf_{Q_3}u\leq \inf_{Q_{1/4}}u\leq 1$. Hence Lemma 5.14 implies
\[
|\{u>\tfrac{P}{2}\}\cap Q_1|\leq C\left(\tfrac{P}{2}\right)^{-\varepsilon}.
\]

We choose $r$ such that $r^{n/2} \ge C(P/2)^{-\varepsilon}$ and $r \le 1/4$. Then we have $Q_r(x_0) \subset Q_1$
and

\[
(1)\qquad \frac{1}{|Q_r(x_0)|}|\{u > P/2\}\cap Q_r(x_0)| \le \frac{1}{2}.
\]

Next, we show that for $\theta > 1$, with $\theta - 1$ small, $u \ge \theta P$ at some point in
$Q_{4\sqrt{n}r}(x_0)$. We prove it by contradiction. Suppose $u \le \theta P$ in $Q_{4\sqrt{n}r}(x_0)$. Consider
the transformation

\[
x = x_0 + ry \qquad \text{for any } y \in Q_{4\sqrt{n}} \text{ and } x \in Q_{4\sqrt{n}r}(x_0),
\]

and the function

\[
\tilde u(y) = \frac{\theta P - u(x)}{(\theta - 1)P}.
\]

Obviously $\tilde u \ge 0$ in $B_{2\sqrt{n}}$ and $\tilde u(0)=1$, hence $\inf_{Q_3}\tilde u \le 1$. It is easy to check that
$\tilde u \in \mathcal{S}^+(\lambda,\Lambda,\tilde f)$ in $B_{2\sqrt{n}}$ with $\|\tilde f\|_{L^n(B_{2\sqrt{n}})} \le \varepsilon_0$. In fact, we have

\[
\tilde f(y) = -\frac{r^2}{(\theta-1)P}f(x) \qquad \text{for any } y \in B_{2\sqrt{n}},
\]

and hence

\[
\|\tilde f\|_{L^n(B_{2\sqrt{n}})} \le \frac{r}{(\theta-1)P}\|f\|_{L^n(B_{2\sqrt{n}r}(x_0))} \le \varepsilon_0,
\]

if we choose $P$ such that $r \le (\theta - 1)P$. Hence, we may apply Lemma 5.13 to $\tilde u$.
Note that $u(x) \le P/2$ if and only if $\tilde u(y) \ge \frac{\theta - 1/2}{\theta - 1}$ and that $\frac{\theta - 1/2}{\theta - 1}$ is large if $\theta$
is close to 1. So we obtain

\[
\frac{1}{|Q_r(x_0)|}|\{u \le P/2\}\cap Q_r(x_0)|
= |\{\tilde u \ge \tfrac{\theta - 1/2}{\theta - 1}\}\cap Q_1| \le C\left(\frac{\theta - 1/2}{\theta - 1}\right)^{-\varepsilon} < \frac{1}{2},
\]

if $\theta$ is chosen close to 1. This contradicts (1).
Hence, we conclude that there exists a $\theta > 1$ such that if

\[
u(x_0)=P \qquad \text{for an } x_0 \in B_{1/4},
\]

then

\[
u(x_1) \ge \theta P \qquad \text{for an } x_1 \in Q_{4\sqrt{n}r}(x_0) \subset B_{2nr}(x_0),
\]

provided

\[
C(n,\lambda,\Lambda)P^{- \varepsilon/n} \le r \le (\theta - 1)P.
\]

where $\theta$ and $C$ are positive constants depending only on $n,\lambda$ and $\Lambda$. We choose $P$
such that $P \ge \left(\frac{C}{\theta-1}\right)^{\frac{n}{n+\varepsilon}}$ and then take $r = CP^{-\varepsilon/n}$.
Now, we iterate the above result to get a sequence $\{x_k\}$ such that for any
$k = 1,2,\cdots,$

\[
u(x_k) \ge \theta^k P \qquad \text{for an } x_k \in B_{2nr_k}(x_{k-1}),
\]

where

\[
r_k = C(\theta^{k-1}P)^{-\varepsilon/n} = C\theta^{-(k-1)\varepsilon/n}P^{-\varepsilon/n}.
\]


In order to have $\{x_k\}\in B_{1/2}$, we need $\sum 2nr_k<1/4$. Hence, we choose $M_0$ such
that
\[
M_0^{\varepsilon/n}\ge 8nC\sum_{k=1}^{\infty}\theta^{-(k-1)\varepsilon/n}
\quad\text{and}\quad
M_0\ge \left(\frac{C}{\theta-1}\right)^{\!n/(n+\varepsilon)},
\]
and then take $P>M_0$. This finishes the proof.

In the rest of this section, we prove a technical lemma concerning the second
order derivatives of functions in $\mathcal{S}(\lambda,\Lambda,f)$. Such a result will be needed in the
discussion of $W^{2,p}$-estimates. First, we introduce some terminology.
Let $\Omega$ be a bounded domain and $u$ be a continuous function in $\Omega$. We define
for $M>0$
\[
\Gm{M}{u,\Omega}=\{x_0\in\Omega;\ \text{there exists an affine function }L\text{ such that}
\]
\[
L(x)-\frac{M}{2}|x-x_0|^2\le u(x)\text{ for }x\in\Omega\text{ with equality at }x_0\},
\]
\[
\Gp{M}{u,\Omega}=\{x_0\in\Omega;\ \text{there exists an affine function }L\text{ such that}
\]
\[
L(x)+\frac{M}{2}|x-x_0|^2\ge u(x)\text{ for }x\in\Omega\text{ with equality at }x_0\},
\]
\[
\G{M}{u,\Omega}=\Gp{M}{u,\Omega}\cap\Gm{M}{u,\Omega}.
\]

We also define
\[
\Am{M}{u,\Omega}=\Omega\setminus\Gm{M}{u,\Omega},
\qquad
\Ap{M}{u,\Omega}=\Omega\setminus\Gp{M}{u,\Omega},
\qquad
\A{M}{u,\Omega}=\Omega\setminus\G{M}{u,\Omega}.
\]

In other words, $\Gm{M}{u,\Omega}$ (resp. $\Gp{M}{u,\Omega}$) consists of points where there is a
concave (resp. convex) paraboloid of opening $M$ touching $u$ from below (resp.
above). Intuitively, $|\A{M}{u,\Omega}|$ behaves like the distribution function of $D^2u$. Hence
for integrability of $D^2u$, we need to study the decay of $|\A{M}{u,\Omega}|$.

\begin{lemma}\label{hanlin:ch05-lem-5-15}\dcref{Lemma 5.15}\group{ch05-viscosity-solutions}\source{han-lin-lecture-notes:page-0111,han-lin-lecture-notes:page-0112,han-lin-lecture-notes:page-0113,han-lin-lecture-notes:page-0114}\uses{hanlin:ch05-lem-5-9,hanlin:ch05-lem-5-16} \textit{Suppose that $\Omega$ is a bounded domain with $B_{6\sqrt{n}}\subset\Omega$ and that $u$
belongs to $S^+(\lambda,\Lambda,f)$ in $B_{6\sqrt{n}}$ for some $f\in C(B_{6\sqrt{n}})$. Then there exist positive
constants $\delta_0$, $\mu$ and $C$, depending only on $n$, $\lambda$ and $\Lambda$, such that, if $|u|\le 1$ in $\Omega$ and
$\|f\|_{L^n(B_{6\sqrt{n}})}\le \delta_0$, there holds}
\[
|\A{t}{u,\Omega}\cap \Qone|\le Ct^{-\mu}\quad\text{for any }t>0.
\]

\textit{If, in addition, $u\in S(\lambda,\Lambda,f)$ in $B_{6\sqrt{n}}$, there holds}
\[
|\A{t}{u,\Omega}\cap \Qone|\le Ct^{-\mu}\quad\text{for any }t>0.
\]
\end{lemma}

In the proof of Lemma 5.15, we need the maximal functions of local integrable
functions. For any $g\in L^1_{\mathrm{loc}}(\mathbb{R}^n)$, we define
\[
m(g)(x)=\sup_{r>0}\frac{1}{|Q_r(x)|}\int_{Q_r(x)}|g|.
\]

The maximal operator $m$ is of weak type $(1,1)$ and of strong type $(p,p)$ for $1<p\leq$
$\infty$, i.e.,

\[
|\{x\in\mathbb{R}^n;m(g)(x)\geq t\}|\leq \frac{c_1(n)}{t}\|g\|_{L^1(\mathbb{R}^n)} \qquad \text{for any } t>0
\]
\[
\|m(g)\|_{L^p(\mathbb{R}^n)}\leq c_2(n,p)\|g\|_{L^p(\mathbb{R}^n)} \qquad \text{for } 1<p\leq \infty.
\]

Now we begin to prove Lemma 5.15. The following result plays an important role. It asserts that, if $u$ has a tangent paraboloid with opening $1$ from below somewhere in $Q_3$, then the set where $u$ has a tangent paraboloid from below with opening $M$ in $Q_1$ is large. Compare it with Lemma 5.13.

\begin{lemma}\label{hanlin:ch05-lem-5-16}\dcref{Lemma 5.16}\group{ch05-viscosity-solutions}\source{han-lin-lecture-notes:page-0112}\uses{hanlin:ch05-lem-5-13,hanlin:ch05-thm-5-8} \textit{Suppose that $\Omega$ is a bounded domain with $B_{6\sqrt{n}}\subset \Omega$ and that $u$ belongs to $\mathcal{S}^+(\lambda,\Lambda,f)$ in $B_{6\sqrt{n}}$ for some $f\in C(B_{6\sqrt{n}})$. Then there exist constants $0<\sigma<1$, $\delta_0>0$ and $M>1$, depending only on $n$, $\lambda$ and $\Lambda$, such that, if $\|f\|_{L^n(B_{6\sqrt{n}})}\leq \delta_0$ and $G_1^-(u,\Omega)\cap Q_3
\neq \phi$, then}

\[
|G_M^-(u,\Omega)\cap Q_1|\geq 1-\sigma.
\]
\end{lemma}

\textsc{Proof.} Since $G_1^-(u,\Omega)\cap Q_3
\neq \phi$, there is an affine function $L_1$ such that
\[
v\geq P_1 \qquad \text{in } \Omega \text{ with equality at some point in } Q_3,
\]
where
\[
v(x)=\frac{u(x)}{2n}+L_1(x) \qquad \text{and} \qquad P_1(x)=1-\frac{|x|^2}{4n}.
\]
This implies $v\geq 0$ in $B_{2\sqrt{n}}$ and $\inf_{Q_3}v\leq 1$. Then as in the proof of Lemma 5.13, for $w=v+g$, where $g$ is the function constructed in Lemma 5.13, we have
\[
|\{w=\Gamma_w\}\cap Q_1|\geq 1-\sigma,
\]
for some $\sigma\in(0,1)$ if $\delta_0$ is chosen small. Now we need to prove for some $M>1$
\[
\{w=\Gamma_w\}\cap Q_1\subset G_M^-(u,\Omega)\cap Q_1.
\]
Let $x_0\in \{w=\Gamma_w\}\cap Q_1$ and take an affine function $L_2$ with $L_2<0$ on $\partial B_{2\sqrt{n}}$ and
\[
L_2\leq \Gamma_w\leq v+g \qquad \text{in } B_{2\sqrt{n}} \text{ with equality at } x_0.
\]
It follows
\[
(1)\qquad P_2\leq L_2-g\leq v \qquad \text{in } B_{2\sqrt{n}} \text{ with equality at } x_0,
\]
for a concave paraboloid $P_2$ of opening $M_0$, a positive constant depending only on
$n$, $\lambda$ and $\Lambda$.

Next, we prove $P_2\leq v$ in $\Omega\setminus B_{2\sqrt{n}}$. Note that $P_2<-g=0=P_1$ on
$\partial B_{2\sqrt{n}}$ and that $P_2(x_0)=v(x_0)\geq P_1(x_0)$ with $x_0\in Q_1\subset B_{2\sqrt{n}}$. If we take
$M_0>1/(2n)$, then $\{P_2-P_1\geq 0\}$ is convex. We conclude $P_2-P_1<0$ in $\mathbb{R}^n\setminus B_{2\sqrt{n}}$.
Hence, we have $P_2\leq P_1\leq v$ in $\Omega\setminus B_{2\sqrt{n}}$. By (1) and the definition of $v$, we get
\[
x_0\in G_{2nM_0}(u,\Omega)\cap Q_1 \text{ with } 2nM_0>1.
\]
\hfill$\square$

\textsc{Proof of Lemma 5.15.} Recall $B_{6\sqrt{n}}\subset \Omega$, $u\in \mathcal{S}^+(\lambda,\Lambda,f)$ in $B_{6\sqrt{n}}$ and
\[
(1)\qquad |u|_{L^\infty(\Omega)}\leq 1,\qquad \|f\|_{L^n(B_{6\sqrt{n}})}\leq \delta_0.
\]
We prove that there exist constants $M>1$ and $0<\gamma<1$, depending only on $n$, $\lambda$
and $\Lambda$, such that
\[
|A_{M^k}^-(u,\Omega)\cap Q_1|\leq \gamma^k \qquad \text{for any } k=0,1,\cdots .
\]


\noindent\textit{Step 1.} There exist constants $M>1$ and $0<\sigma<1$ such that
\begin{equation}
\tag{2}
\left|G_M^-(u,\Omega)\cap Q_1\right|\ge 1-\sigma.
\end{equation}
It is easy to see that $\|u\|_{L^\infty(\Omega)}\le 1$ implies
\[
G_c^-(u,\Omega)\cap Q_3
\neq \varnothing,
\]
for some constant $c$ depending only on $n$. We now apply Lemma 5.16 to $u/c$ to get (2). By a simple adjustment, we assume that $\delta_0$, $M$ and $\sigma$ in Step 1 are the same as those in Lemma 5.15.


\noindent\textit{Step 2.} We extend $f$ by zero outside $B_{6\sqrt n}$ and set for $k=0,1,\cdots$,
\[
A=A_{M^{-(k+1)}}^-(u,\Omega)\cap Q_1,
\]
\[
B=\left(A_{M^{-k}}^-(u,\Omega)\cap Q_1\right)\cup\{x\in Q_1;m(f^n)(x)\ge (c_1M^k)^n\},
\]
for some $c_1>0$ to be determined. Then we claim
\[
|A|\le \sigma|B|,
\]
where $M>1$ and $0<\sigma<1$ are as before. Recall that $m(f^n)$ denotes the maximal function of $f^n$.

We prove it by Lemma 5.9. It is easy to see $|A|\le \sigma$ since we have $\left|G_{M^{k+1}}^-(u,\Omega)\cap Q_1\right|\ge \left|G_M^-(u,\Omega)\cap Q_1\right|\ge 1-\sigma$ by Step 1. Next, we claim that, if $Q=Q_r(x_0)$ is a cube in $Q_1$ such that
\begin{equation}
\tag{3}
\left|A_{M^{-(k+1)}}^-(u,\Omega)\cap Q\right|=\left|A\cap Q\right|>\sigma|Q|,
\end{equation}
then $\widetilde Q\cap Q_1\subset B$ for $\widetilde Q=Q_{3r}(x_0)$. We prove it by contradiction. If not, we take an $\widetilde x$ such that
\[
\widetilde x\in G_{M^{-k}}^-(u,\Omega)\cap \widetilde Q,
\]
and
\[
\sup_{r>0}\frac{1}{|Q_r(\widetilde x)|}\int_{Q_r(\widetilde x)}|f|^n\le (c_1M^k)^n.
\]
Consider the transformation
\[
x=x_0+ry\qquad\text{for any }y\in Q_1\text{ and }x\in Q=Q_r(x_0),
\]
and the function
\[
\widetilde u(y)=\frac{1}{r^2M^k}u(x).
\]
It is easy to check that $B_{6\sqrt n}\subset \widetilde\Omega$, the image of $\Omega$ under the transformation above, and that $\widetilde u\in \mathcal S^+(\lambda,\Lambda,\widetilde f)$ in $B_{6\sqrt n}$ with
\[
\widetilde f(y)=\frac{1}{M^k}f(x)\qquad\text{for any }y\in B_{6\sqrt n}.
\]
By the choice of $\widetilde x$, we have
\[
G_1^-(\widetilde u,\widetilde\Omega)\cap Q_3
\neq \varnothing.
\]
Since $B_{6\sqrt n r}(x_0)\subset Q_{15\sqrt n r}(\widetilde x)$, there holds
\[
\|\widetilde f\|_{L^n(B_{6\sqrt n})}\le \frac{1}{rM^k}\|f\|_{L^n(Q_{15\sqrt n r}(\widetilde x))}\le c(n)c_1\le \delta_0,
\]
if we take $c_1$ small enough, depending only on $n$, $\lambda$ and $\Lambda$.

\[
\text{Hence }\tilde u \text{ satisfies the assumption of Lemma 5.16 with }\Omega\text{ replaced by }\tilde\Omega.\text{ We}
\]
\[
\text{apply Lemma 5.16 to }\tilde u\text{ to get}
\]
\[
\lvert G_M^-(\tilde u,\tilde\Omega)\cap Q_1\rvert \geq 1-\sigma,
\]
\[
\text{or}
\]
\[
\lvert G_{M^{k+1}}^-(u,\Omega)\cap Q\rvert > (1-\sigma)\lvert Q\rvert.
\]

This contradicts (3). We are in a position to apply Lemma 5.9.

\emph{Step 3.} We finish the proof of Lemma 5.15. Define for $k = 0,1,\cdots,$
\[
\alpha_k = \lvert A_{M^k}(u,\Omega)\cap Q_1\rvert,
\]
\[
\beta_k = \lvert\{x\in Q_1 : m(f^n)(x)\geq (c_1M^k)^n\}\rvert.
\]
Then Step 2 implies $\alpha_{k+1}\leq \sigma(\alpha_k+\beta_k)$ for any $k = 0,1,\cdots$. Hence by an iteration,
we have
\[
\alpha_k \leq \sigma^k + \sum_{i=0}^{k-1}\sigma^{k-i}\beta_i.
\]

Since $\lVert f^n\rVert_{L^1}\leq \delta_0^n$ and the maximal operator is of the weak type $(1,1)$, we conclude
\[
\beta_k \leq c(n)\delta_0^n(c_1M^k)^{-n} = CM^{-nk},
\]
where $C$ is a positive constant depending only on $n$, $\lambda$ and $\Lambda$. This implies
\[
\sum_{i=0}^{k-1}\sigma^{k-i}\beta_i \leq C\sum_{i=0}^{k-1}\sigma^{k-i}M^{-ni}\leq Ck\gamma_0^k,
\]
with $\gamma_0 = \max\{\sigma,M^{-n}\}<1$. Therefore, we obtain for $k$ large
\[
\alpha_k \leq \sigma^k + Ck\gamma_0^k \leq (1+Ck)\gamma_0^k \leq \gamma^k,
\]
for some constant $\gamma\in(0,1)$, depending only on $n$, $\lambda$ and $\Lambda$. This finishes the
proof.

\[
\square
\]

\begin{remark}\label{hanlin:ch05-rem-5-17}\dcref{Remark 5.17}\group{ch05-viscosity-solutions}\source{han-lin-lecture-notes:page-0114}\uses{hanlin:ch05-lem-5-15} The polynomial decay of the function
\[
\mu(t) = \lvert A_t(u,\Omega)\cap Q_1\rvert
\]
for $u\in \mathcal S(\lambda,\Lambda,f)$ implies that $D^2u$ is $L^p$-integrable in $Q_1$ for small $p>0$, depending
only on $n$, $\lambda$ and $\Lambda$. In order to show the $L^p$-integrability for large $p$, we need to
speed up the convergence in the proof of Lemma 5.15. We will discuss $W^{2,p}$-
estimates in Section 5.4.
\end{remark}

\begin{center}
\textbf{5.3. Schauder Estimates}
\end{center}

In this section, we prove the Schauder estimates for viscosity solutions.
Throughout this section, we always assume that $a_{ij}\in C(B_1)$ satisfies
\[
\lambda |\xi|^2 \leq a_{ij}(x)\xi_i\xi_j \leq \Lambda |\xi|^2 \qquad \text{for any }x\in B_1\text{ and any }\xi\in \mathbb R^n,
\]
for some positive constants $\lambda$ and $\Lambda$ and that $f$ is a continuous function in $B_1$.
The following approximation result plays an important role in the discussion
of the regularity theory.

\begin{lemma}\label{hanlin:ch05-lem-5-18}\dcref{Lemma 5.18}\group{ch05-viscosity-solutions}\source{han-lin-lecture-notes:page-0115}\uses{hanlin:ch05-thm-5-8,hanlin:ch05-cor-5-11,hanlin:lemma-holder-boundary} Suppose $u \in C(B_1)$ is a viscosity solution of
\[
a_{ij}D_{ij}u=f \qquad \text{in } B_1,
\]
with $|u| \leq 1$ in $B_1$. Assume for some $0<\varepsilon<1/16$,
\[
\|a_{ij}-a_{ij}(0)\|_{L^n(B_{3/4})}\leq \varepsilon.
\]
Then there exists a function $h \in C(\overline{B_{3/4}})$, with $a_{ij}(0)D_{ij}h=0$ in $B_{3/4}$ and $|h|\leq 1$ in $B_{3/4}$, such that
\[
\|u-h\|_{L^\infty(B_{1/2})}\leq C\{\varepsilon^\gamma+\|f\|_{L^n(B_1)}\},
\]
where $C$ and $\gamma\in(0,1)$ are constants depending only on $n$, $\lambda$ and $\Lambda$.
\end{lemma}

Proof. Solve for $h\in C(\overline{B_{3/4}})\cap C^\infty(B_{3/4})$ such that
\[
a_{ij}(0)D_{ij}h=0 \qquad \text{in } B_{3/4}
\]
\[
h=u \qquad \text{on } \partial B_{3/4}.
\]
The maximum principle implies $|h|\leq 1$ in $B_{3/4}$. Note that $u$ belongs to $S(\lambda,\Lambda,f)$ in $B_1$. Corollary 5.11 implies $u\in C^\alpha(\overline{B_{3/4}})$ for some $\alpha\in(0,1)$ depending only on $n$, $\lambda$ and $\Lambda$, with the estimate
\[
\|u\|_{C^\alpha(\overline{B_{3/4}})}\leq C\{1+\|f\|_{L^n(B_1)}\},
\]
where $C$ is a positive constant depending only on $n$, $\lambda$ and $\Lambda$. By Lemma 1.36 in Chapter 1, we have
\[
\|h\|_{C^{\alpha/2}(\overline{B_{3/4}})}\leq C\|u\|_{C^\alpha(\overline{B_{3/4}})}\leq C\{1+\|f\|_{L^n(B_1)}\}.
\]
Since $u-h=0$ on $\partial B_{3/4}$, we get for any $0<\delta<1/4$
\[
\tag{1}
\|u-h\|_{L^\infty(\partial B_{3/4-\delta})}\leq C\delta^{\alpha/2}\{1+\|f\|_{L^n(B_1)}\}.
\]
We claim for any $0<\delta<1$
\[
\tag{2}
\|D^2 h\|_{L^\infty(B_{3/4-\delta})}\leq C\delta^{\alpha/2-2}\{1+\|f\|_{L^n(B_1)}\}.
\]
In fact, for any $x_0\in B_{3/4-\delta}$, we apply interior $C^2$-estimate to $h-h(x_1)$ in $B_\delta(x_0)\subset B_{3/4}$ for some $x_1\in \partial B_\delta(x_0)$ and obtain
\[
|D^2 h(x_0)|\leq C\delta^{-2}\sup_{B_\delta(x_0)}|h-h(x_1)|\leq C\delta^{-2}\delta^{\alpha/2}\{1+\|f\|_{L^n(B_1)}\}.
\]
Note that $u-h$ is a viscosity solution of
\[
a_{ij}D_{ij}(u-h)=f-(a_{ij}-a_{ij}(0))D_{ij}h\equiv F \qquad \text{in } B_{3/4}.
\]
By Theorem 5.8, Alexandroff maximum principle, we have with (1) and (2)
\begin{align*}
\|u-h\|_{L^\infty(B_{3/4-\delta})}
&\leq \|u-h\|_{L^\infty(\partial B_{3/4-\delta})}+C\|F\|_{L^n(B_{3/4-\delta})}\\
&\leq \|u-h\|_{L^\infty(\partial B_{3/4-\delta})}+C\|D^2 h\|_{L^\infty(B_{3/4-\delta})}\|a_{ij}-a_{ij}(0)\|_{L^n(B_{3/4})}+C\|f\|_{L^n(B_1)}\\
&\leq C(\delta^{\alpha/2}+\delta^{\alpha/2-2}\varepsilon)\{1+\|f\|_{L^n(B_1)}\}+C\|f\|_{L^n(B_1)}.
\end{align*}
Take $\delta=\varepsilon^{1/2}<1/4$ and then $\gamma=\alpha/4$. This finishes the proof. $\square$

For the next result, we need to introduce the following concept.

\begin{definition}\label{hanlin:ch05-def-5-19}\dcref{Definition 5.19}\group{ch05-viscosity-solutions}\source{han-lin-lecture-notes:page-0116} A function $g$ is Hölder continuous at $0$ with exponent $\alpha$ in
$L^n$-sense if
\[
[g]_{C^\alpha_{L^n}}(0)\equiv \sup_{0<r<1}\frac{1}{r^\alpha}\left(\frac{1}{|B_r|}\int_{B_r}|g-g(0)|^n\right)^{1/n}<\infty.
\]
\end{definition}

Now we state the Schauder estimates.

\begin{theorem}\label{hanlin:ch05-thm-5-20}\dcref{Theorem 5.20}\group{ch05-viscosity-solutions}\source{han-lin-lecture-notes:page-0116,han-lin-lecture-notes:page-0117,han-lin-lecture-notes:page-0118}\uses{hanlin:ch05-def-5-19,hanlin:ch05-lem-5-18} Suppose $u\in C(B_1)$ is a viscosity solution of
\[
a_{ij}D_{ij}u=f \qquad \text{in } B_1.
\]
Assume $\{a_{ij}\}$ is Hölder continuous at $0$ with exponent $\alpha$ in $L^n$-sense for some
$\alpha\in(0,1)$. If $f$ is Hölder continuous at $0$ with exponent $\alpha$ in $L^n$-sense, then $u$ is
$C^{2,\alpha}$ at $0$. Moreover, there exists a polynomial $P$ of degree $2$ such that
\[
|u-P|_{L^\infty(B_r(0))}\le Cr^{2+\alpha} \qquad \text{for any } 0<r<1,
\]
\[
|P(0)|+|DP(0)|+|D^2P(0)|\le C_*,
\]
and
\[
C_*\le C\{|u|_{L^\infty(B_1)}+|f(0)|+[f]_{C^\alpha_{L^n}}(0)\},
\]
where $C$ is a positive constant depending only on $n,\lambda,\Lambda,\alpha$ and $[a_{ij}]_{C^\alpha_{L^n}}(0)$.
\end{theorem}

Proof. First we assume $f(0)=0$. For this, we consider $v=u-b_{ij}x_i x_j f(0)/2$
for a constant matrix $\{b_{ij}\}$ such that $a_{ij}(0)b_{ij}=1$. By scaling, we also assume that
$[a_{ij}]_{C^\alpha_{L^n}}(0)$ is small. Next, by considering $\delta>0$
\[
\frac{u}{|u|_{L^\infty(B_1)}+\frac{1}{\delta}[f]_{C^\alpha_{L^n}}(0)},
\]
we may assume $|u|_{L^\infty(B_1)}\le 1$ and $[f]_{C^\alpha_{L^n}}(0)\le \delta$.

In the following, we prove that there is a constant $\delta>0$, depending only on
$n,\lambda,\Lambda$ and $\alpha$, such that, if $u\in C(B_1)$ is a viscosity solution of
\[
a_{ij}D_{ij}u=f \qquad \text{in } B_1,
\]
with
\[
|u|_{L^\infty(B_1)}\le 1,\qquad [a_{ij}]_{C^\alpha_{L^n}}(0)\le \delta,
\]
\[
\left(\frac{1}{|B_r|}\int_{B_r}|f|^n\right)^{1/n}\le \delta r^\alpha \qquad \text{for any } 0<r<1,
\]
then there exists a polynomial $P$ of degree $2$ such that
\[
(1)\qquad |u-P|_{L^\infty(B_r(0))}\le Cr^{2+\alpha} \qquad \text{for any } 0<r<1,
\]
and
\[
(2)\qquad |P(0)|+|DP(0)|+|D^2P(0)|\le C,
\]
where $C$ is a positive constant depending only on $n,\lambda,\Lambda$ and $\alpha$.

We claim that there exist a $\mu\in(0,1)$, depending only on $n,\lambda,\Lambda$ and $\alpha$, and a
sequence of polynomials of degree $2$ of the form
\[
P_k(x)=a_k+b_k\cdot x+\frac12 x^tC_kx,
\]

such that for any $k = 0,1,2,\cdots,$
\[
a_{ij}(0)D_{ij}P_k = 0,
\]
\begin{equation}\tag{3}
\lvert u - P_k\rvert_{L^\infty(B_{\mu^k})} \le \mu^{k(2+\alpha)},
\end{equation}
and
\begin{equation}\tag{4}
\lvert a_k - a_{k-1}\rvert + \mu^{k-1}\lvert b_k - b_{k-1}\rvert + \mu^{2(k-1)}\lvert C_k - C_{k-1}\rvert \le C\mu^{(k-1)(2+\alpha)},
\end{equation}
where $P_0 = P_{-1} = 0$, and $C$ is a positive constant depending only on $n,\lambda,\Lambda$ and $\alpha$.

We first prove that Theorem 5.20 follows from (3) and (4). It is easy to see that $a_k,b_k$ and $C_k$ converge and that the limiting polynomial
\[
p(x) = a_\infty + b_\infty \cdot x + \frac{1}{2}x^t C_\infty x
\]
satisfies
\[
\lvert P_k(x) - p(x)\rvert \le C\bigl(\lvert x\rvert^2\mu^{\alpha k} + \lvert x\rvert\mu^{(\alpha+1)k} + \mu^{(\alpha+2)k}\bigr) \le C\mu^{(2+\alpha)k},
\]
for any $\lvert x\rvert \le \mu^k$. Hence, we have for $\lvert x\rvert \le \mu^k$
\[
\lvert u(x) - p(x)\rvert \le \lvert u(x)-P_k(x)\rvert + \lvert P_k(x)-p(x)\rvert \le C\mu^{(2+\alpha)k},
\]
and hence
\[
\lvert u(x)-p(x)\rvert \le C\lvert x\rvert^{2+\alpha} \qquad \text{for any } x \in B_1.
\]

Now we prove (3) and (4). Clearly (3) and (4) hold for $k = 0$. We assume they hold for $k = 0,1,2,\cdots,l$ and proceed to prove for $k = l+1$. Consider the function
\[
\tilde{u}(y) = \frac{1}{\mu^{l(2+\alpha)}}\bigl(u - P_l\bigr)(\mu^l y) \qquad \text{for any } y \in B_1.
\]
Then $\tilde{u} \in C(B_1)$ is a viscosity solution of
\[
\tilde{a}_{ij}D_{ij}\tilde{u} = \tilde{f} \qquad \text{in } B_1,
\]
with
\[
\tilde{a}_{ij}(y) = \frac{1}{\mu^{l\alpha}}a_{ij}(\mu^l y),
\]
\[
\tilde{f}(y) = \frac{1}{\mu^{l\alpha}}\bigl(f(\mu^l y) - a_{ij}(\mu^l y)D_{ij}P_k\bigr).
\]

Now we check that $\tilde{u}$ satisfies the assumptions of Lemma 5.18. For this, we note
\[
\lVert \tilde{a}_{ij} - \tilde{a}_{ij}(0)\rVert_{L^n(B_1)} \le \frac{1}{\mu^{l\alpha}}\lVert a_{ij} - a_{ij}(0)\rVert_{L^n(B_{\mu^l})} \le [a_{ij}]_{C^\alpha_{L^n}(0)} \le \delta,
\]
and
\[
\lVert \tilde{f}\rVert_{L^n(B_1)} \le \frac{1}{\mu^{l\alpha}}\lVert f\rVert_{L^n(B_{\mu^l})} + \frac{1}{\mu^{l\alpha}}\sup \lvert D^2P_l\rvert \,\lVert a_{ij} - a_{ij}(0)\rVert_{L^n(B_{\mu^l})} \le \delta + C\delta,
\]
where we used
\[
\lvert D^2P_l\rvert \le \sum_{k=1}^l \lvert D^2P_k - D^2P_{k-1}\rvert \le \sum_{k=1}^l \mu^{(k-1)\alpha} \le C.
\]
Hence we take $\varepsilon = C(n,\lambda,\Lambda)\delta$ as in Lemma 5.18. Then by Lemma 5.18, there exists a function $h \in C(\bar{B}_{3/4})$ with $\tilde{a}_{ij}(0)D_{ij}h = 0$ in $B_{3/4}$ and $\lvert h\rvert \le 1$ in $B_{3/4}$ such that
\[
\lVert \tilde{u} - h\rVert_{L^\infty(B_{1/2})} \le C(\varepsilon^\gamma + \varepsilon) \le 2C\varepsilon^\gamma.
\]

Write $\tilde{P}(y) = h(0) + Dh(0) + y^t D^2h(0)y/2$. Then by interior estimates for $h$ we have
\[
\lvert \tilde{u} - \tilde{P}\rvert_{L^\infty(B_\mu)} \le \lvert \tilde{u} - h\rvert_{L^\infty(B_\mu)} + \lvert h - \tilde{P}\rvert_{L^\infty(B_\mu)} \le 2C\varepsilon^\gamma + C\mu^3 \le \mu^{2+\alpha},
\]
by choosing $\mu$ small and then $\varepsilon$ small accordingly. Rescaling back, we have
\[
\lvert u(x) - P_l(x) - \mu^{l(2+\alpha)}\tilde{P}(\mu^{-l}x)\rvert \le \mu^{(l+1)(2+\alpha)} \qquad \text{for any } x \in B_{\mu^{l+1}}.
\]
This implies (3) for $k = l + 1$, if we define
\[
P_{k+1}(x) = P_k(x) + \mu^{l(2+\alpha)}\tilde{P}(\mu^{-l}x).
\]
The estimate (4) follows easily. $\square$

To finish this section, we state a Cordes-Nirenberg type estimate. The proof is similar to that of Theorem 5.20.

\begin{theorem}\label{hanlin:ch05-thm-5-21}\dcref{Theorem 5.21}\group{ch05-viscosity-solutions}\source{han-lin-lecture-notes:page-0118}\uses{hanlin:ch05-thm-5-20} \emph{Suppose $u \in C(B_1)$ is a viscosity solution of}
\[
a_{ij}D_{ij}u = f \qquad \text{in } B_1.
\]
\emph{Then for any $\alpha \in (0,1)$, there exists a $\theta > 0$, depending only on $n,\lambda,\Lambda$ and $\alpha$, such that if}
\[
\left(\frac{1}{\lvert B_r\rvert}\int_{B_r}\lvert a_{ij} - a_{ij}(0)\rvert^n\right)^{\frac{1}{n}} \le \theta \qquad \text{for any } 0 < r \le 1,
\]
\emph{then $u$ is $C^{1,\alpha}$ at $0$; that is, there exists an affine function $L$ such that}
\[
\lvert u - L\rvert_{L^\infty(B_r(0))} \le C_*r^{1+\alpha} \qquad \text{for any } 0 < r < 1,
\]
\[
\lvert L(0)\rvert + \lvert DL(0)\rvert \le C_*,
\]
\emph{and}
\[
C_* \le C\left\{\lvert u\rvert_{L^\infty(B_1)} + \sup_{0<r<1} r^{1-\alpha}\left(\frac{1}{\lvert B_r\rvert}\int_{B_r}\lvert f\rvert^n\right)^{\frac{1}{n}}\right\},
\]
\emph{where $C$ is a positive constant depending only on $n,\lambda,\Lambda$ and $\alpha$.}
\end{theorem}

\textbf{5.4. $W^{2,p}$ Estimates}

In this section, we prove $W^{2,p}$-estimates for viscosity solutions. Throughout this section, we always assume that $a_{ij} \in C(B_1)$ satisfies
\[
\lambda\lvert\xi\rvert^2 \le a_{ij}(x)\xi_i\xi_j \le \Lambda\lvert\xi\rvert^2 \qquad \text{for any } x \in B_1 \text{ and any } \xi \in \mathbb{R}^n,
\]
for some positive constants $\lambda$ and $\Lambda$ and that $f$ is a continuous function in $B_1$.

The main result in this section is the following theorem.

\begin{theorem}\label{hanlin:ch05-thm-5-22}\dcref{Theorem 5.22}\group{ch05-viscosity-solutions}\source{han-lin-lecture-notes:page-0118,han-lin-lecture-notes:page-0119}\uses{hanlin:ch05-lem-5-23} \emph{Suppose $u \in C(B_1)$ is a viscosity solution of}
\[
a_{ij}D_{ij}u = f \qquad \text{in } B_1.
\]
\emph{Then for any $p \in (n,\infty)$, there exists an $\epsilon > 0$, depending only on $n,\lambda,\Lambda$ and $p$, such that if}
\[
\left(\frac{1}{\lvert B_r(x_0)\rvert}\int_{B_r(x_0)}\lvert a_{ij} - a_{ij}(x_0)\rvert^n\right)^{\frac{1}{n}} \le \epsilon \qquad \text{for any } B_r(x_0) \subset B_1,
\]

then $u \in W^{2,p}_{\mathrm{loc}}(B_1)$. Moreover,

\[
\|u\|_{W^{2,p}(B_{1/2})} \leq C\{|u|_{L^\infty(B_1)} + \|f\|_{L^p(B_1)}\},
\]

where $C$ is a positive constant depending only on $n$, $\lambda$, $\Lambda$ and $p$.
\end{theorem}

As before, we prove the following result instead.

\begin{lemma}\label{hanlin:ch05-lem-5-23}\dcref{Lemma 5.23}\group{ch05-viscosity-solutions}\source{han-lin-lecture-notes:page-0119,han-lin-lecture-notes:page-0121,han-lin-lecture-notes:page-0122}\uses{hanlin:ch05-lem-5-15,hanlin:ch05-lem-5-24} Suppose $u \in C(B_{8\sqrt{n}})$ is a viscosity solution of
\[
a_{ij}D_{ij}u = f \qquad \text{in } B_{8\sqrt{n}}.
\]

Then for any $p \in (n,\infty)$, there exist positive constants $\varepsilon$ and $C$, depending only on
$n$, $\lambda$, $\Lambda$ and $p$, such that if
\[
\|u\|_{L^\infty(B_{8\sqrt{n}})} \leq 1, \qquad \|f\|_{L^p(B_{8\sqrt{n}})} \leq \varepsilon,
\]
and
\[
\left(\frac{1}{|B_r(x_0)|}\int_{B_r(x_0)} |a_{ij}-a_{ij}(x_0)|^n\right)^{1/n} \leq \varepsilon
\qquad \text{for any } B_r(x_0) \subset B_{8\sqrt{n}},
\]
then $u \in W^{2,p}(B_1)$ and $\|u\|_{W^{2,p}(B_1)} \leq C$.
\end{lemma}

We first describe the strategy of the proof. Let $\Omega$ be a bounded domain and $u$
be a continuous function in $\Omega$. As in Section 5.2, we define for $M>0$
\[
G_M(u,\Omega) = \{x_0 \in \Omega;\ \text{there exists an affine function } L \text{ such that}
\]
\[
L(x) - \frac{M}{2}|x-x_0|^2 \leq u(x) \leq L(x) + \frac{M}{2}|x-x_0|^2
\]
\[
\text{for } x \in \Omega \text{ with equality at } x_0\},
\]
\[
A_M(u,\Omega) = \Omega \setminus G_M(u,\Omega).
\]

We consider the function
\[
\theta(x) = \theta(u,\Omega)(x) = \inf\{M; x \in G_M(u,\Omega)\} \in [0,\infty]
\qquad \text{for } x \in \Omega.
\]

It is straightforward to verify that for $p \in (1,\infty]$ the condition $\theta \in L^p(\Omega)$ implies
$D^2u \in L^p(\Omega)$ and
\[
\|D^2u\|_{L^p(\Omega)} \leq 2\|\theta\|_{L^p(\Omega)}.
\]

In order to study the integrability of the function $\theta$, we discuss its distribution
function
\[
\mu_\theta(t) = |\{x \in \Omega; \theta(x) > t\}| \qquad \text{for any } t>0.
\]

It is clear that
\[
\mu_\theta(t) \leq |A_t(u,\Omega)| \qquad \text{for any } t>0.
\]

Hence we need to study the decay of $|A_t(u,\Omega)|$.

\begin{lemma}\label{hanlin:ch05-lem-5-24}\dcref{Lemma 5.24}\group{ch05-viscosity-solutions}\source{han-lin-lecture-notes:page-0119,han-lin-lecture-notes:page-0120,han-lin-lecture-notes:page-0121}\uses{hanlin:ch05-lem-5-15,hanlin:ch05-lem-5-18} Suppose that $\Omega$ is a bounded domain with $B_{8\sqrt{n}} \subset \Omega$ and that
$u \in C(\Omega)$ is a viscosity solution of
\[
a_{ij}D_{ij}u = f \qquad \text{in } B_{8\sqrt{n}}.
\]

Then for any $\varepsilon_0 \in (0,1)$, there exist an $M>1$, depending only on $n$, $\lambda$ and $\Lambda$, and
an $\varepsilon \in (0,1)$, depending only on $n$, $\lambda$, $\Lambda$ and $\varepsilon_0$, such that if
\[
(1) \qquad \|f\|_{L^n(B_{8\sqrt{n}})} \leq \varepsilon,\qquad \|a_{ij}-a_{ij}(0)\|_{L^n(B_{7\sqrt{n}})} \leq \varepsilon,
\]

and

(2)

\[
G_1(u,\Omega)\cap Q_3\neq \emptyset,
\]

then

\[
|G_M(u,\Omega)\cap Q_1|\ge 1-\varepsilon_0.
\]
\end{lemma}


\noindent\textbf{Proof.} Let $x_1\in G_1(u,\Omega)\cap Q_3$. Then, there exists an affine function $L$ such
that
\[
-\frac12|x-x_1|^2\le u(x)-L(x)\le \frac12|x-x_1|^2 \qquad \text{in }\Omega.
\]

By considering $(u-L)/c$ instead of $u$, for a constant $c>1$ large enough, depending
only on $n$, we may assume

\[
(3)\qquad |u|\le 1 \qquad \text{in }B_{6\sqrt n},
\]

and hence

\[
(4)\qquad -|x|^2\le u(x)\le |x|^2 \qquad \text{for any }x\in \Omega\setminus B_{6\sqrt n}.
\]

Solve for $h\in C(\bar B_{7\sqrt n})\cap C^\infty(B_{7\sqrt n})$ such that
\[
a_{ij}(0)D_{ij}h=0 \qquad \text{in }B_{7\sqrt n}
\]
\[
h=u \qquad \text{on }\partial B_{7\sqrt n}.
\]

Then Lemma 5.18 implies

\[
(5)\qquad |u-h|_{L^\infty(B_{6\sqrt n})}\le C\{\varepsilon^\gamma+\|f\|_{L^n(B_{6\sqrt n})}\},
\]

and

\[
(6)\qquad \|h\|_{C^2(B_{6\sqrt n})}\le C,
\]

where $C>0$ and $\gamma\in(0,1)$, as in Lemma 5.18, depending only on $n,\lambda$ and $\Lambda$.
Consider $h|_{\bar B_{6\sqrt n}}$. Extend $h$ outside $\bar B_{6\sqrt n}$ continuously such that $h=u$ in $\Omega\setminus B_{7\sqrt n}$
and $|u-h|_{L^\infty(\Omega)}=|u-h|_{L^\infty(B_{6\sqrt n})}$. Note $|h|\le 1$ in $\Omega$. It follows that $|u-h|_{L^\infty(\Omega)}\le
2$ and hence with (4)

\[
-2-|x|^2\le h(x)\le 2+|x|^2 \qquad \text{for any }x\in \Omega\setminus \bar B_{6\sqrt n}.
\]

Then there exists an $N>1$, depending only on $n,\lambda$ and $\Lambda$, such that

\[
(7)\qquad Q_1\subset G_N(h,\Omega).
\]

Consider

\[
w=\frac{\min\{1,\delta_0\}}{2C\varepsilon^\gamma}(u-h),
\]

where $\delta_0$ is the constant in Lemma 5.15, and $C$ and $\gamma$ are constants in (5) and (6).
It is easy to check that $w$ satisfies the hypothesis of Lemma 5.15 in $\Omega$. We now
apply Lemma 5.15 to get

\[
|A_t(w,\Omega)\cap Q_1|\le Ct^{-\mu} \qquad \text{for any }t>0.
\]

Therefore, we have

\[
|A_s(u-h,\Omega)\cap Q_1|\le C\varepsilon^\gamma s^{-\mu} \qquad \text{for any }s>0.
\]

It follows that

\[
|G_N(u-h,\Omega)\cap Q_1|\ge 1-C_1\varepsilon^\gamma\ge 1-\varepsilon_0,
\]

if we choose $\varepsilon\in(0,1)$ small, depending only on $n$, $\lambda$, $\Lambda$ and $\varepsilon_0$. With (7) we get

\[
|G_{2N}(u,\Omega)\cap Q_1|\ge 1-\varepsilon_0.
\]

This finishes the proof.

\hfill$\square$

\begin{remark}\label{hanlin:ch05-rem-5-25}\dcref{Remark 5.25}\group{ch05-viscosity-solutions}\source{han-lin-lecture-notes:page-0121}\uses{hanlin:ch05-lem-5-24} In fact, we proved Lemma 5.24 with the assumption (2) replaced
by (3).
\end{remark}

\textbf{Proof of Lemma 5.23.} \textit{Step 1.} For any $\varepsilon_0\in(0,1)$, there exist an $M>1$,
depending only on $n$, $\lambda$ and $\Lambda$, and an $\varepsilon\in(0,1)$, depending only on $n$, $\lambda$, $\Lambda$ and $\varepsilon_0$,
such that, under the assumptions of Lemma 5.23, there holds

\[
(1)\qquad |G_M(u,B_{8\sqrt n})\cap Q_1|\ge 1-\varepsilon_0.
\]

We remark that $M$ does not depend on $\varepsilon_0$. In fact, we have $|u|\le 1\le |x|^2$ in
$B_{8\sqrt n}\setminus B_{6\sqrt n}$. We may apply Lemma 5.24 to get (1) with $\Omega=B_{8\sqrt n}$ (see Remark
5.25).

\textit{Step 2.} We set, for $k=0,1,\cdots$,

\[
A=A_{M^{k+1}}(u,B_{8\sqrt n})\cap Q_1,
\]
\[
B=\bigl(A_{M^k}(u,B_{8\sqrt n})\cap Q_1\bigr)\cup \{x\in Q_1;m(f^n)(x)\ge (c_1M^k)^n\},
\]

for some $c_1>0$ to be determined, depending only on $n$, $\lambda$, $\Lambda$ and $\varepsilon_0$. Then there
holds

\[
|A|\le \varepsilon_0|B|.
\]

The proof is identical to that of Lemma 5.15.

\textit{Step 3.} We finish the proof of Lemma 5.24. We take $\varepsilon_0$ such that

\[
\varepsilon_0M^p=\frac12,
\]

where $M$, depending only on $n$, $\lambda$ and $\Lambda$, is as in Step 1. Hence the constants $\varepsilon$ and
$c_1$ depend only on $n$, $\lambda$, $\Lambda$ and $p$. Define for $k=0,1,\cdots$,

\[
\alpha_k=|A_{M^k}(u,B_{8\sqrt n})\cap Q_1|,
\]

\[
\beta_k=|\{x\in Q_1;m(f^n)(x)\ge (c_1M^k)^n\}|.
\]

Then Step 2 implies $\alpha_{k+1}\le \varepsilon_0(\alpha_k+\beta_k)$ for any $k=0,1,\cdots$. Hence, by an iteration
we have

\[
\alpha_k\le \varepsilon_0^k+\sum_{i=1}^{k-1}\varepsilon_0^{k-i}\beta_i.
\]

Since $f^n\in L^{p/n}$ and the maximal operator is of strong type $(p,p)$, we conclude
that $m(f^n)\in L^{p/n}$ and

\[
\|m(f^n)\|_{L^{p/n}}\le C\|f\|_{L^p}^n\le C.
\]

Then the definition of $\beta_k$ implies

\[
\sum_{k\ge 0}M^{pk}\beta_k\le C.
\]

As before, we set

\[
\theta(x)=\theta(u,B_{1/2})(x)=\inf\{M;x\in G_M(u,B_{1/2})\}\in[0,\infty]\qquad \text{for }x\in B_{1/2},
\]

and

\[
\mu_\theta(t)=|\{x\in B_{1/2};\ \theta(x)>t\}|\qquad \text{for any }t>0.
\]

The proof will be finished if we show

\[
\|\theta\|_{L^p(B_{1/2})}\le C.
\]

It is clear that

\[
\mu_\theta(t)\le |A_t(u,B_{1/2})|\le |A_t(u,B_{8\sqrt n})\cap Q_1|\qquad \text{for any }t>0.
\]

It suffices to prove, with the definition of $\alpha_k$, that

\[
\sum_{k\ge 1}M^{pk}\alpha_k\le C.
\]

In fact, we have

\[
\sum_{k\ge 1}M^{pk}\alpha_k\le \sum_{k\ge 1}(\varepsilon_0M^p)^k+\sum_{k\ge 1}\sum_{i=0}^{k-1}\varepsilon_0^{k-i}M^{p(k-i)}M^{pi}\beta_i
\]

\[
\le \sum_{k\ge 1}2^{-k}+\Bigl(\sum_{i\ge 0}M^{pi}\beta_i\Bigr)\Bigl(\sum_{j\ge 1}2^{-j}\Bigr)\le C.
\]

This finishes the proof.

\hfill$\square$
